% Generated draft: requires mathematical and source review.
\documentclass[11pt]{article}
\usepackage[margin=25mm]{geometry}
\usepackage{fontspec}
\setmainfont{Latin Modern Roman}
\usepackage{amsmath,amssymb,mathtools,unicode-math}
\setmathfont{Latin Modern Math}
\usepackage{enumitem,needspace,xurl,hyperref}
\setlength{\parindent}{0pt}
\setlength{\parskip}{5pt}
\setlength{\emergencystretch}{3em}
\newcommand{\N}{\mathbb{N}}
\newcommand{\Z}{\mathbb{Z}}
\newcommand{\Q}{\mathbb{Q}}
\newcommand{\R}{\mathbb{R}}
\newcommand{\C}{\mathbb{C}}
\newcommand{\F}{\mathbb{F}}
\newcommand{\A}{\mathbb{A}}
\newcommand{\PP}{\mathbb{P}}
\newcommand{\E}{\mathbb{E}}
\newcommand{\eps}{\varepsilon}
\newcommand{\dd}{\,\mathrm{d}}
\begin{document}
\section*{UnsolvedMath ID: 701690}
\textbf{Source problem code:} catalogue entry
\subsection{Definitions and mathematical statement}
Let \(G\) be a group with identity \(e\), and let \(A\subseteq G\) be a symmetric generating set, meaning
\[
A=A^{-1}:=\{a^{-1}:a\in A\},\qquad \langle A\rangle=G.
\]
The literal finite-list question in the supplied record is whether, for every such pair \((G,A)\), there exist a positive integer \(n\) and elements \(a_1,\ldots,a_n\in A\) such that the map
\[
\{1,\ldots,n\}\longrightarrow G,\qquad
k\longmapsto p_k:=a_1a_2\cdots a_k
\]
is a bijection. Thus the partial products must be pairwise distinct and exhaust \(G\); the generators \(a_i\) themselves need not be distinct.

\textbf{Unresolved domain qualification.} The source does not state that \(G\) is finite. A finite list exhausting \(G\) entails finiteness and \(n=|G|\). If finite groups are the intended domain, the universal quantifier above must explicitly be restricted to finite groups; this restriction is not an express hypothesis of the supplied record.

\textbf{Boundary and graph qualifications.} The displayed list begins at \(p_1=a_1\), not at the empty product \(p_0=e\). The identity must occur exactly once among \(p_1,\ldots,p_n\), but the source does not require \(p_n=e\) or specify a closing edge. Neither \(e\notin A\) nor \(A\ne\varnothing\) nor \(|G|>1\) is imposed. The introductory assertion concerning the Lov\'asz conjecture uses the word ``Hamiltonian'' without specifying a path or a cycle. Its intended relationship to the displayed partial-product question remains unresolved; no additional closing condition is imposed here.
\subsection{Short English statement}
Given a symmetric generating set of a group, can one choose a finite sequence of its generators whose successive products enumerate every group element exactly once? The source leaves finiteness unstated and does not specify a Hamiltonian path or cycle convention. This is a statement-only draft requiring human mathematical review.
\subsection{Sources}
\begin{itemize}
\item Supplied local UnsolvedMath registry record: permanent ID \(701690\), exact source problem code \texttt{catalogue entry}, title ``Kourovka Notebook Problem 20.32.'' The registry labels the problem \texttt{open}; this is a supplied status, not an independently verified mathematical or literature assessment.
\item The supplied background attributes Problem 20.32 to \emph{The Kourovka Notebook (Unsolved Problems in Group Theory)}, 20th issue (2022), and the 21st issue edition (2026), identifying \texttt{arXiv:1401.0300v48}: \url{https://arxiv.org/abs/1401.0300v48}. The supplied collection description names E. I. Khukhro and V. D. Mazurov as editors. This bibliographic attribution has not been independently checked.
\item The proposer fields and background name B. Green, with proposed year 2022. The statement separately begins with the attribution ``(E. Rapaport Strasser).'' Both attributions are retained; their respective roles are not explained by the supplied record.
\item The source invokes the Lov\'asz conjecture in its introductory motivation. No separate citation or precise path-versus-cycle formulation is supplied for that assertion.
\end{itemize}
\end{document}
